\section*{अध्याय \ref{ch:g9:acids-bases-ph} --- अम्ल, क्षारक और pH}

\begin{solution}{exo:g9:acids-bases-ph:1}
\omterm{def:g9:acids-bases-ph:ph}{pH} 3: अम्लीय। \omterm{def:g9:acids-bases-ph:ph}{pH} 7: उदासीन। \omterm{def:g9:acids-bases-ph:ph}{pH} 11: क्षारकीय।
\end{solution}

\begin{solution}{exo:g9:acids-bases-ph:2}
अम्लीय: \ce{H+} \omterm{def:g9:ions:ion}{आयन}। क्षारकीय: \ce{OH-} \omterm{def:g9:ions:ion}{आयन}।
\end{solution}

\begin{solution}{exo:g9:acids-bases-ph:3}
कागज़ी नींबू का रस (2), कॉफ़ी (5), शुद्ध पानी (7), समुद्र का पानी (8.1), घरेलू
अमोनिया (12)।
\end{solution}

\begin{solution}{exo:g9:acids-bases-ph:4}
साफ़, सूखी तश्तरी पर रखी पट्टी को साफ़ काँच की छड़ से लाई गई विलयन की एक बूँद
से छुआया जाता है; रंग की तुलना तुरंत
चार्ट से की जाती है।
\end{solution}

\begin{solution}{exo:g9:acids-bases-ph:5}
दस-गुना तनुकरण के बाद लगभग 5; सौ-गुना के बाद लगभग 6।
\end{solution}

\begin{solution}{exo:g9:acids-bases-ph:6}
नहीं। तनु करने से \omterm{def:g9:acids-bases-ph:ph}{pH} 7 के क़रीब आता है, पर उसके पार कभी नहीं जाता: पानी ख़ुद
उदासीन है, और उसे मिलाने से \ce{OH-} \omterm{def:g9:ions:ion}{आयन} \ce{H+}
\omterm{def:g9:ions:ion}{आयनों} से अधिक नहीं हो सकते।
\end{solution}

\begin{solution}{exo:g9:acids-bases-ph:7}
मिलाने पर ऊष्मा निकलती है। सांद्र अम्ल पर डाला गया पानी तुरंत उबलकर अम्ल को
बाहर उछाल सकता है; पानी में धीरे-धीरे डाला गया अम्ल तुरंत बड़े आयतन में
फैल जाता है।
\end{solution}

\begin{solution}{exo:g9:acids-bases-ph:8}
संक्षारक: त्वचा और आँखों को जलाता है, धातुओं पर आक्रमण करता है। चश्मा और दस्ताने
पहनो; कोई भी छींटा ढेर सारे पानी से धो डालो।
\end{solution}

\begin{solution}{exo:g9:acids-bases-ph:9}
तीन तनुकरण (1 से 2, 2 से 3, 3 से 4): $10 \times 10 \times 10 =
1000$ गुना।
\end{solution}

\begin{solution}{exo:g9:acids-bases-ph:10}
मिल्क ऑफ़ मैग्नीशिया क्षारकीय है (\omterm{def:g9:acids-bases-ph:ph}{pH} लगभग 10.5): वह पेट के अतिरिक्त अम्ल के
कुछ भाग के साथ अभिक्रिया करता है।
\end{solution}

\begin{solution}{exo:g9:acids-bases-ph:11}
नहीं: 8.1 के \omterm{def:g9:acids-bases-ph:ph}{pH} के साथ वह अब भी थोड़ा क्षारकीय है। वह कम \omterm{def:g9:acids-bases-ph:ph}{pH} की ओर, यानी कम
क्षारकीय होने की ओर, बढ़ रहा है, क्योंकि \omterm{def:g4:air-a-mixture-of-gases:air}{वायु} की कार्बन डाइऑक्साइड उसमें
\omterm{def:g2:mixing-and-dissolving:dissolve}{घुलती} है।
\end{solution}

\begin{solution}{exo:g9:acids-bases-ph:12}
\omterm{def:g9:acids-bases-ph:ph-paper}{pH पत्र} को रंग-चार्ट से आँख द्वारा मिलाकर पढ़ा जाता है, लगभग एक इकाई तक; \omterm{def:g9:acids-bases-ph:ph-paper}{pH
मीटर} एक या दो दशमलव तक की संख्या देता है। मीटर अधिक परिशुद्ध है; दोनों परिणाम
पत्र की परिशुद्धता के भीतर मेल खाते हैं।
\end{solution}

\begin{solution}{pb:g9:acids-bases-ph:1}
\textbf{1.} अम्लीय: उसमें ज़्यादातर \ce{H+(aq)} और \ce{Cl-(aq)}
\omterm{def:g9:ions:ion}{आयन} हैं।

\textbf{2.} GHS05, संक्षारक: अम्ल त्वचा, आँखों और धातुओं पर आक्रमण करता है।

\textbf{3.} वह आँख से मिलाए जाने वाले रंग के बजाय एक परिशुद्ध संख्या
देता है।

\textbf{4.} दस गुना: \omterm{def:g9:acids-bases-ph:ph}{pH} की एक इकाई दस-गुना तनुकरण के संगत है।

\textbf{5.} \qty{10}{L}, \omterm{def:g9:acids-bases-ph:ph}{pH} 3 का।

\textbf{6.} तीन: \omterm{def:g9:acids-bases-ph:ph}{pH} 2 से 3, 3 से 4, 4 से 5।

\textbf{7.} $1 \times 10 \times 10 \times 10 = \qty{1000}{L}$।

\textbf{8.} नहीं: तनु करने से \omterm{def:g9:acids-bases-ph:ph}{pH} केवल 7 के क़रीब आता है।

\textbf{9.} \ce{H+ + OH- -> H2O}।

\textbf{10.} अम्ल एक विशाल आयतन में फैलने के बजाय नष्ट होकर पानी बन जाता है;
थोड़ा-सा क्षारक ही चाहिए, हज़ार लीटर पानी
नहीं।

\textbf{11.} पानी में घुला खाने का सोडा (या साबुन वाला पानी)।

\textbf{12.} $1000 - 1 = \qty{999}{L}$ पानी।
\end{solution}
